% Source: https://artofproblemsolving.com/community/c3h2533664_not_a_valid_solution
% Stable identity: aops-topic-2533664-q01
\textbf{(a)} Among $1,\dots,6$, the pairs that are not relatively prime are the pairs of even digits ($2,4$; $2,6$; $4,6$) and the pair $3,6$. So we need: no two even digits adjacent, and $3$ and $6$ not adjacent.

The even digits occupy three pairwise non-adjacent positions among $1,\dots,6$. The possible position sets are $\{1,3,5\}$, $\{2,4,6\}$, $\{1,3,6\}$ and $\{1,4,6\}$, i.e. the patterns
\[eoeoeo,\qquad oeoeoe,\qquad eoeooe,\qquad eooeoe.\]
Each pattern can be filled in $3!\cdot3!=36$ ways, giving $4\cdot36=144$ numbers with no two even digits adjacent.

Now subtract those in which $3$ and $6$ are adjacent. If $6$ is in a position with $j$ odd neighbours, then $3$ can go in $j$ ways and the other two even and two odd digits in $2!\cdot2!=4$ ways. Adding the numbers of odd neighbours of the even positions:
\begin{itemize}
\item $eoeoeo$: $1+2+2=5$, giving $20$; similarly $oeoeoe$ gives $20$;
\item $eoeooe$: $1+2+1=4$, giving $16$; similarly $eooeoe$ gives $16$.
\end{itemize}
So $20+20+16+16=72$ numbers have $3$ next to $6$, and the answer is $144-72=\boxed{72}$.

\textbf{(b)} The quicker solution makes two errors that happen to cancel.
\begin{itemize}
\item It only counts the alternating patterns $eoeoeo$ and $oeoeoe$, and misses the patterns $eoeooe$ and $eooeoe$, which also keep the even digits apart. These contain $2\cdot(36-16)=40$ valid numbers, for example $216534$.
\item It never checks the pair $3,6$, so it counts alternating numbers such as $214365$, where $3$ and $6$ are adjacent. Of the $72$ alternating numbers, $40$ are invalid and only $32$ are valid.
\end{itemize}
So it leaves out $40$ valid numbers and includes $40$ invalid ones; getting $72$ is a coincidence, not a proof.
